\documentclass[a4paper,12pt]{article}
\usepackage[spanish,activeacute]{babel}
\usepackage[latin1]{inputenc}

%% Para las columnas
\usepackage{multicol}


\usepackage{amssymb, amsmath}
\usepackage{latexsym}
\usepackage{multicol}


%% Para numeraci?n autom?tica
\newcounter{nroproblema}
\setcounter{nroproblema}{1}
\newcounter{nroapartado}
\setcounter{nroapartado}{1}
\newcommand{\n}{%
    \vspace{.1in}\noindent{\bf \arabic{nroproblema}.\ }%
    \addtocounter{nroproblema}{1}%
    \setcounter{nroapartado}{1}%
    }%
\newcommand{\apartado}{\vspace*{0.2cm}\hspace*{0.2cm}%
   {\bf \alph{nroapartado})}
   \addtocounter{nroapartado}{1}%
   }

%% m?genes

\setlength{\unitlength}{1mm} \addtolength{\voffset}{-15mm}
\addtolength{\textheight}{30mm} \addtolength{\hoffset}{-15mm}
\addtolength{\textwidth}{30mm}


\begin{document}

\begin{center}{\bf {\Large \'ALGEBRA I. HOJA 9}}\end{center}

\n Hallar los autovalores de las siguientes matrices, y diagonalizar cuando
sea posible (es decir, hallar una matriz inversible $P$ tal que $A = PDP^{-1}$).

a) $\left( \begin{matrix}3 & 4 \\
- 2 & - 3\end{matrix} \right)$.

b) $\left( \begin{matrix}2 &  -1 \\
1 & 0\end{matrix} \right)$.

c) $\left( \begin{matrix} - 2 &  9 & -6 \\
1 & -2 & 0\\
 3 & -9 & 5 \end{matrix} \right)$.
 
d) $\left( \begin{matrix} 2 &  -1  & -1 \\
0 & 3 & 2\\
-1 & 1 & 2 \end{matrix} \right)$.

e) $\left( \begin{matrix} 1 &  1  & 0 \\
 0 & 1 & 1\\
0 & 0 & 1 \end{matrix} \right)$.

f) $\left( \begin{matrix} -10 &  6  & 3 \\
-26 & 16 & 8\\
16 & -10 & -5 \end{matrix} \right)$.


g) $\left( \begin{matrix} 0 &  -6  & -16 \\
 0 & 17 & 45\\
0 & -6 & -16 \end{matrix} \right)$.
 
 Probar que las matrices en f) y g) son similares.
 
 \n  Diagonalizar $A=\left( \begin{matrix}1 & 4 \\
2 & 3\end{matrix} \right)$. Calcular $A^5$.

Respuesta parcial: los autovalores son 5 y -1.

\n Sea $a\in\mathbb{R}$, $a\ne 0$. Demostrar que
$\left( \begin{matrix}1 & a \\
0 & 1\end{matrix} \right)$ no es diagonalizable.

\n  Probar que si  $a\in\mathbb{R}$
no es un m�ltiplo entero de $2\pi$, entonces la matriz $\left( \begin{matrix}\cos a & -\sin a \\
\sin a & \cos a\end{matrix} \right)$ no tiene autovectores.

\n  Probar que si  $a\in\mathbb{R}$, entonces la matriz $A := \left( \begin{matrix}\cos a & \sin a \\
\sin a & - \cos a\end{matrix} \right)$  tiene (al menos) un autovector $v$ asociado al autovalor
$1$. Probar que si $w$ es perpendicular a $v$, entonces $A w = - w$.



\end{document}